\section{进动、TiP与声子语言的边界}\label{sec:3-4}

在第\ref{sec:3-2}节与第\ref{sec:3-3}节中，我们建立了纯三轴转子及单准粒子耦合转子绕稳定主轴摆动运动的局域展开模型，推导了摆动特征频率在横向与纵向几何构型下的演化机制。在非轴对称原子核的转动谱学中，围绕低激发伴带的动力学本质长期存在着术语内涵与模型机制的多重竞争：一方面，唯象声子图景常将激发带直接赋予“单声子摆动（$n_w=1$）”或“双声子摆动（$n_w=2$）”标签，但在有限自旋的介观核体系中，简谐声子近似的有效性边界、波函数的空间弥散度以及内禀自由度的重取向效应需要经受严格的理论审视；另一方面，在正常形变核区（如 $A \approx 130$ 核区），伴带与基带之间的能级间距随自旋演化的特征究竟应当被理解为几何各向异性主导的横向摆动，还是归因于三维空间转轴倾斜进动（Tilted Precession, TiP）、多准粒子组态交叉、普通旋称伙伴带分裂或 $\gamma$ 软核的大振幅形状涨落，成为了低能核结构谱学的前沿争论焦点。

本节旨在系统阐明轴对称高 $K$ 能带进动的微观机制，严格解耦固定自旋局域旋转系振动频率与固定内禀投影跨自旋实验室能量间距；深入阐释倾斜进动（TiP）模型的粒子—转子取向几何，将有限自旋量子波函数的节点与概率分布同经典相空间轨道清晰区分；进而划定简谐声子语言在有限核中的适用界限与多重替代机制的条件竞争；最后确立结合多维电磁跃迁网络进行实验鉴别与联合判据的三轴谱学分析准则。

\subsection{轴对称高 $K$ 能带进动与转动态能量间距}\label{sec:3-4-1}

在探讨非轴对称核取向动力学之前，必须首先建立轴对称变形核中高 $K$ 态的严格物理基准。对于理想的轴对称形状转子，绕其对称轴（取为3轴）的形状取向是一个冗余自由度坐标；唯象模型以约束条件 $R_3 = 0$ 表达体系不存在独立的轴向集体形状转动。需要强调，这一形状转动约束并不意味着任意轴对称微观体系的轴向推转响应或有效转动惯量必然处处为零，微观单粒子自由度依然可以沿对称轴承载显著的轨道与自旋取向。

设核内未配对价准粒子的总内禀角动量算符为 $\hat{\mathbf{J}}_{\text{intr}} = \sum_i (\hat{\mathbf{l}}_i + \hat{\mathbf{s}}_i)$。定义无量纲角动量算符 $\hat{\boldsymbol{\mathcal{I}}} \equiv \hat{\mathbf{I}}/\hbar$ 与 $\hat{\mathbf{j}} \equiv \hat{\mathbf{J}}_{\text{intr}}/\hbar$，其在内禀对称轴（3轴）上的无量纲内禀投影量子数为 $K$（$K \gg 1$）。核芯集体转动角动量 $\hat{\mathbf{R}}$ 具有作用量物理量纲，其垂直于对称轴的分量满足 $\hat{\mathbf{R}}_\perp = \hbar (\hat{\boldsymbol{\mathcal{I}}}_\perp - \hat{\mathbf{j}}_\perp)$。设体系在垂直于对称轴方向的集体主转动惯量为 $\mathcal{J}_\perp > 0$。集体转动动能仅由垂直于对称轴的集体角动量支撑：
\begin{equation}\label{eq:3-4-high-K-H}
\hat{H}_{\text{coll}} = \frac{\hat{\mathbf{R}}_\perp^2}{2\mathcal{J}_\perp} = \frac{\hbar^2}{2\mathcal{J}_\perp} \left( \hat{\boldsymbol{\mathcal{I}}}_\perp - \hat{\mathbf{j}}_\perp \right)^2.
\end{equation}

在固定内禀组态及确定 $K$ 的强耦合绝热近似下，忽略与其他内禀组态的非对角科里奥利混合，并将固定内禀组态的常数项贡献吸收到内禀能 $E_{\text{intr}}$，该强耦合模型给出的实验室系高 $K$ 转动能谱为：
\begin{equation}\label{eq:3-4-high-K-spectrum}
E_K(I) - E_{\text{intr}} = \frac{\hbar^2}{2\mathcal{J}_\perp} \left[ I(I+1) - K^2 \right] \quad (I = K, K+1, K+2, \dots).
\end{equation}
在此模型中，固定 $K$ 与常数 $\mathcal{J}_\perp$ 构成了明确的计算条件；并非仅由单个对角基矢矩阵元 $\langle K | \hat{j}_\pm | K \rangle = 0$ 即可推知能够忽略所有态混合。在允许的旋称与宇称选择定则下，该模型的理论带头自旋满足 $I_{\text{bandhead}} = K$。在实际在束 $\gamma$ 谱学观测中，由于分支退激阈值、退耦衰变以及低自旋微观能级扰动，实验测得的最低带能级不必然严格等于内禀量子数 $K$。

考察该转动带内相邻自旋能级在实验室系中的跃迁能量间距：
\begin{equation}\label{eq:3-4-lab-spacing}
\Delta E_{\text{lab}}(I \to I-1) \equiv E_K(I) - E_K(I-1) = \frac{\hbar^2 I}{\mathcal{J}_\perp}.
\end{equation}
式~\eqref{eq:3-4-lab-spacing} 在所考查的两个能级均为该模型允许的带内态时成立；若将此能差对应到某条实际光子，还须另检验其选择定则和跃迁矩阵元。若 $K+1$ 与 $K$ 均为允许能级，带头首跨能级跃迁间距为：
\begin{equation}\label{eq:3-4-first-lab-spacing}
\Delta E_{\text{lab}}(K+1 \to K) = \frac{\hbar^2 (K+1)}{\mathcal{J}_\perp}.
\end{equation}

在此必须严格厘清局域共转参考系特征频率与实验室跨自旋能级差分的本质物理区别：
\begin{enumerate}
\item \textbf{局域频率与跨自旋能差的比较对象}：第\ref{sec:3-2}节在固定自旋 $I$ 的纯转子模型中展开横向取向自由度，所得局域频率刻画该自旋子空间中的小振幅激发；式~\eqref{eq:3-2-limit2-nonzero}给出了其中一个惯量简并极限。式~\eqref{eq:3-4-lab-spacing}则比较固定内禀 $K$ 的转动带内不同自旋 $I$ 与 $I-1$ 的实验室能级。两者之间的关系需要指定哈密顿量及参考系，下面的轴对称 RPA 提供一个具体例子；
\item \textbf{轴对称 RPA 参考系关系与进动角频率}：Shimizu、Matsuzaki 与 Matsuyanagi（2005）~\cite{shimizu-2005-high-k-precession} 系统探讨了微观准粒子随机相位近似（RPA）摆动形式理论的轴对称极限。为与该文的坐标选择一致，下段将同一对称轴暂记为 $x$ 轴，横向记为 $y,z$，不直接沿用前段的数字轴号。当形变在该转轴周围恢复轴对称性（$\mathcal{J}_y = \mathcal{J}_z \equiv \mathcal{J}_\perp$）且内禀投影满足 $K = \langle \Phi | \hat{J}_x | \Phi \rangle / \hbar$ 时，对于物理的 $\Delta K = +1$ 激发模，实验室进动角频率与均匀推转系角频率满足精确的参考系运动学变换关系：
\begin{equation}\label{eq:3-4-shimizu-connection}
\omega_{\text{wob}} = \omega_{\text{prec}} - \omega_{\text{rot}}, \qquad \omega_{\text{prec}} = \frac{\hbar K}{\mathcal{J}_\perp^{\text{eff}}(\omega_{\text{prec}})}.
\end{equation}
需要明确的是，式~\eqref{eq:3-4-shimizu-connection} 中的参考系关系来自轴对称小振幅 RPA 运动方程，并不要求预先取绝热极限；有效转动惯量 $\mathcal{J}_\perp^{\text{eff}}(\omega_{\text{prec}})$ 显式依赖于模式固有频率，后一式需自洽求解。在该文讨论的进一步绝热极限，可将这一频率相关有效惯量近似为相应的静态推转响应。在固定惯量强耦合能谱与 $K \gg 1$、低阶激发的共同适用域内，$\hbar \omega_{\text{prec}} \simeq \hbar^2 K / \mathcal{J}_\perp$ 与实验室带头跃迁间距的主导项渐近趋同，而精确的首能级差仍保留 $\hbar^2(K+1)/\mathcal{J}_\perp$ 的量子修正。
\end{enumerate}

\subsection{倾斜进动（TiP）模型的取向动力学描述}\label{sec:3-4-2}

针对正常形变核区（如 $A \approx 130$ 区）中伴带能量间隔随自旋反常演化的物理机制，Lawrie、Shirinda 与 Petrache（2020）基于粒子—转子模型（PRM）建立了倾斜进动（Tilted Precession, TiP）唯象取向动力学图景~\cite{lawrie-2020-tilted-precession-wobbling}。

在粒子—转子模型中，总哈密顿量同时包含核芯转动与内禀自由度。对纯转子，或将内禀角动量冻结为常矢量等简化经典实现，可将固定角动量球面与相应二次能量曲面的交线用作取向运动的几何示意。对内禀角动量也随时间演化的一般 PRM，这个低维交线图不能代替完整的耦合运动。
\begin{itemize}
\item \textbf{准粒子取向与核芯转动的动力学竞争}：奇 $A$ 核中侵入的高 $j$ 准粒子（如缺中子钡镧铈区中的 $\pi h_{11/2}$ 轨道）受到强各向异性四极驱动势的作用，倾向于顺排在能量最有利的几何主轴（粒子型倾向于沿短轴 $s$，空穴型倾向于沿长轴 $l$）；而三轴形变核芯由于集体流体力学效应，倾向于绕其有效转动惯量最大的主轴（如无旋流模型下的中间轴 $m$）转动。价核子优势顺排取向与核芯惯量优势转轴的夹角竞争，导致系统的经典稳定转动取向轴脱离单一几何主轴，向主轴所张成的平面（如 $s$-$m$ 平面）发生空间倾斜；
\item \textbf{TiP 能带的物理定义与进动取向}：总角动量矢量围绕这一倾斜的稳定方向做三维空间进动。在 Lawrie 等人讨论的取向序列中，量子态分布可借助倾斜方向与进动幅度作解释，作者将相应三维取向运动描述为 TiP；这不把任意 QTR 激发态等同于唯一的经典进动锥；
\item \textbf{摆动近似的局域性质}：在所讨论的取向模型中，围绕稳定方向作小振幅展开可得到近似谐振模式。横向角动量涨落相对于总自旋应足够小，高阶能量项与算符修正应受控；低阶能差及适当归一的跃迁强度接近谐振标度，可作为这一近似的检验。有限自旋修正、耦合模式或非线性跃迁算符可以改变这些标度，不能把严格线性设为所有取向激发的定义；
\item \textbf{模型机制与唯象判据的客观界限}：准粒子取向与核芯转动的张角竞争是粒子—转子哈密顿量本征求解给出的内部模型机制，并非可以在实验室系中利用探测器沿时间流逝直接追踪的经典质点轨迹。此外，理论计算表明，倾斜进动带与基带之间的能级间距 $\Delta E(I)$ 随自旋演化的趋势（增加、下降或非单调改变）敏感地依赖于具体的单粒子轨道、核芯惯量排序及有效形变参数，能级差下降并非区分 TiP 与其他模式的排他性唯一判据；同理，亦不能脱离微观自洽输入，为主轴三轴度或自旋窗口强加跨核区通用的绝对经验界限。
\end{itemize}

\subsection{有限自旋波函数分布与经典轨迹的几何区分}\label{sec:3-4-3}

在探讨三轴核的取向激发时，极易将经典力学相空间中绘制的连续轨道切线，直观等同于微观核态的实际量子概率分布。在有限自旋的介观原子核中，必须确立严格的量子力学几何区分：
\begin{enumerate}
\item \textbf{经典相轨道的局限性}：在选定的经典哈密顿量与守恒量下，普通周期轨道具有确定的瞬时取向。分离轨道可为非周期极限，更一般的耦合系统也不必给出单一闭合曲线；这些轨迹与量子本征态概率密度是不同的描述；
\item \textbf{一维反射对称性与中心节点条件}：在一维局域坐标 $x$ 上，若系统的有效哈密顿量在空间反射 $x \mapsto -x$ 下保持不变（动能、势及边界条件共同保持该反射），且所考查的本征态为该一维反射下的奇态（$\psi(-x) = -\psi(x)$），则由该反射奇性在原点固定点处得到 $\psi(0) = -\psi(0) \implies \psi(0) = 0$。这一中心节点的出现是反射奇本征态的严格对称性推论，绝不局限于理想的一维简谐势阱——只要势阱保持对称，高阶对称非简谐项（如 $x^4$）自身绝对不会消除中心节点。理想谐振子态的附加特殊性在于其第一激发态还展现出精确的对称双峰结构；
\item \textbf{有限自旋核量子态的节点与弥散特征}：对于处于三维旋转空间中的有限自旋多体核态，角分布或边缘化概率密度是否在某一空间方向严格为零，必须依据完整的多体波函数、角动量投影以及实际定义的离散对称性进行严格评估：
\begin{itemize}
\item 对同一一维反射奇态，中心节点的结论失去前提，可能是相关反射不再是哈密顿量的对称性、所考查态不再具有确定奇性，或所画的是不同的投影/边缘概率。某个科里奥利项或通道耦合是否破坏该反射，须按实际算符逐项检查；仅有各向异性或非简谐项不构成判据；
\item 角向坐标反射、空间三维反演宇称 $\hat{\Pi}$ 以及空间旋转旋称 $\mathcal{R}_x(\pi)$ 属于完全不同的物理对称性操作，绝对不能将它们之间的通道混合混同为同一机制；
\item 在有限自旋态下，量子波函数由离散的非对角基矢相干展开，基态波包在局部稳定二次极小邻域的横向角向展宽满足 $\Delta \theta \sim \Delta \mathcal{I}_\perp / I \sim 1/\sqrt{I}$，该标度关系仅适用于远离软模与分离线的小振幅低阶激发态，不能视为所有自旋区通用的高斯型半高全宽。
\end{itemize}
\item \textbf{实验可观测量的客观映射}：在核物理实验中，探测器记录的是光子能量、时间戳与角度符合计数，并不存在直接对空间量子概率云密度或瞬时经典轨道进行直接成像的测量手段。理论计算求得的量子态密度仅能作为模型解释的微观媒介，对核结构取向动力学的科学判定必须完全建立在可实验验证的跃迁能谱、分支比、多极混合比及约化电磁跃迁几率 $B(E2)$ 和 $B(M1)$ 之上。
\end{enumerate}

\subsection{谐振声子语言的适用界限与替代机制竞争}\label{sec:3-4-4}

基于前述取向动力学与波函数特性的分析，我们可以严格划定简谐声子语言在核谱学中的适用边界与多重竞争机制。

\subsubsection{声子标记的局域有效性与破坏条件}

在微观粒子—转子模型或宏观转子展开中，将激发带标记为 $n_w = 1$ 或 $n_w = 2$ 摆动声子带，在动力学上仅在满足以下局域有效条件时才构成合理的物理近似：
\begin{itemize}
\item \textbf{有效谐振势阱展开条件}：系统在选定的转动取向参考轴邻域内，总有效势能面具有清晰的局域极小，且量子波包的主体被紧密约束在二次型势阱展开的有效作用半径之内，使得高阶非简谐微扰项（如四次方拉伸项与非对角反常耦合）相对于谐振能标 $\hbar \omega_{\text{wob}}$ 属于小量；
\item \textbf{正规模稳定性与微观响应界限}：局域有效恢复力由惯量张量与内禀取向共同支撑。小振幅声子展开并不要求三个主转动惯量必须显著互异（如前文所论证的轴对称极限 $\mathcal{J}_y = \mathcal{J}_z$ 同样存在严格自洽的局域 RPA 声子展开）。其核心物理要求在于，体系在选定自旋邻域内具有稳定且局部隔离的微小振动正规模；在冻结顺排假定下，准粒子优势取向相对核芯惯量轴保持稳定；而在包含微观组态动力学耦合或推转自洽响应的更完备框架中（如第\ref{sec:3-6}节所分析的振动正规模），即便包含频率依赖的自洽极化或动态重排，只要小振幅微扰展开保持收敛且未发生大振幅模式交叠，系统依然能够表现出谐振特征。不能将沿自旋序列的静态参数演化与体系在时间域的瞬态演化混为一谈，亦不能将冻结取向或绝热退耦作为所有谐振声子描述的先验绝对判据。
\end{itemize}

当体系处于低自旋区（$I \sim j$）或遭遇多准粒子带交叉区时，强烈的科里奥利力与准粒子退耦使得势能曲面变得极其平坦甚至发生极值点漂移，非简谐非对角相互作用急剧增强。然而，带交叉或低自旋扰动的出现并不意味着在所有物理情景下都会导致电磁跃迁强度的彻底淬灭或声子物理图像的瞬间瓦解；实验与微观计算应当针对具体的能级拓扑与波函数成分，实事求是地评估声子概念的适用程度及其偏差。

\subsubsection{谱学中的多重竞争机制}

在在束 $\gamma$ 谱学中，当实验测得一条伴带具有密集的 $\Delta I = 1$ 退激并与基带保持较小能量差时，除横向摆动之外，以下替代物理机制构成了强有力的条件竞争解释：
\begin{enumerate}
\item \textbf{普通旋称伙伴带（Signature Partner Bands）}：在单奇 $A$ 形变核中，同一高 $j$ 准粒子轨道（如 $\pi h_{11/2}$）分裂出的两个旋称分支（$\alpha = \pm 1/2$）在适度形变与费米面位置下，其旋称分裂能量差随自旋升高可以表现出平缓缩小或交替振荡的行为，在特定自旋区间内能够模仿出伴带能级差随自旋下降的表象；
\item \textbf{相邻准粒子组态交叉（Band Crossing）}：奇 $A$ 核在高自旋转动过程中，极易经历破对（例如一对高 $j$ 中子对或质子对的顺排）引发的多准粒子组态交叉。在能级交叉区域，由于两组态之间的残余相互作用发生能级排斥与波函数混合，往往在有限自旋跨度内诱发伴带能谱的局部压缩甚至近简并；
\item \textbf{软核芯耦合与不同模型自由度}：对所讨论的一些过渡核，未配对准粒子与 $\gamma$ 软核芯的耦合可在没有深陷刚性三轴极小的条件下给出低激发带及电磁预言；Nomura（2022）的 IBFM 提供了具体实现~\cite{nomura-2022-questioning-wobbling-ibfm}。Lv（2022）的固定形变 QTR 则允许准粒子重新取向，属于另一种竞争实现，不能将其等同于显式的动态 $\gamma$ 振动计算~\cite{lv-2022-evidence-against-wobbling-135pr}。不同模型的自由度、标定输入和未拟合预言应分别比较。
\end{enumerate}

在实际核谱学研究中，必须严格恪守实验数据的测量边界与不确定度传播规律。若相关实验测量仅确立了激发能级纲图与相对分支比，而缺乏能级平均寿命（如反冲距离法 RDM 或多普勒位移衰减法 DSAM）以及 Coulomb 激发等绝对跃迁矩阵元数据，则该数据集无法独立提供带间跃迁的绝对电磁几率 $B(E2)$ 与 $B(M1)$。在数据不完备的情形下，研究必须明确标注测量缺口与物理边界，坚决避免在缺乏直接绝对测量支撑下虚构或武断推断绝对跃迁强度。

\subsection{实验观测量裁决能力与三轴谱学分析准则}\label{sec:3-4-5}

面对上述多重竞争模型的挑战，单一观测量（如孤立的能谱能级差或定性的跃迁存在性）无法对三轴取向动力学做出排他性裁决。现代核结构谱学确立了联合多维实验可观测量进行综合判决的分析准则：
\begin{enumerate}
\item \textbf{能谱间距与激发趋势（$\Delta E(I)$）}：提供动力学模式的初步判据，建立伴带相对于基带在整个可测自旋区间的演化轮廓，用以初筛是否存在旋称反转、带交叉或特征频率分叉；
\item \textbf{方向角关联（DCO 比值）与线偏振度（Linear Polarization）}：用于测定级联退激的跃迁多极度（区分电四极 $E2$ 与磁偶极 $M1$）及多极混合比 $\delta(E2/M1)$。在固定符合门条件与探测几何下，定向核角分布与方向角关联比值 $R_{\text{DCO}}(\delta) = W(\theta_1, \theta_2, \phi; \delta) / W(\theta_2, \theta_1, \phi; \delta)$ 表现为两个关于多极混合比 $\delta$ 的二次式之比（Lange、Kumar 与 Hamilton 1982）~\cite{lange-kumar-hamilton-1982-multipole-admixtures}。求解实验代数方程 $R_{\text{DCO}}(\delta) = R_{\text{exp}}$ 时：
\begin{itemize}
\item 交叉相乘后的方程须区分二次、一次和完全退化的情形，并剔除原分母为零的伪根。在非退化二次情形，判别式决定有两个、一个重合或没有有限实根。纯 $E2$ 极限的 $|\delta|\to\infty$ 是振幅比坐标的端点，不能与有限分母零点混称；具体反演见第\ref{sec:5-7}节；
\item 测量 Compton 散射线偏振度 $P$ 能够提供电性与磁性辐射性质的互补鉴别信息，但其测量的是光子辐射特征，并非能级固有空间宇称的直接实测；在偏振探测灵敏度有限或理论偏振值处于实验误差包络内时，测量 $P$ 不保证能彻底消除多解歧义性。此外，源自同批符合事例重构的 DCO 比值与偏振度数据在统计上并不天然相互独立，综合误差分析中须注意关联效应；
\item 多极混合比 $\delta(E2/M1)$ 严格定义为电四极与磁偶极辐射振幅之比，必须与第\ref{sec:5-9}节采用相同的国际单位制（SI）算符与规范相匹配。若磁偶极跃迁使用传统的磁矩型约化矩阵元 $\mathcal{M}_\mu(M1)$（以核磁子 $\mu_N$ 为单位），则相应的电荷型磁算符定义为 $\mathcal{M}_S(M1) \equiv \mathcal{M}_\mu(M1)/c$（具有电荷乘长度量纲 $e\cdot\mathrm{fm}$）。式~\eqref{eq:3-4-delta-def} 中的 $E_\gamma/(\hbar c)$ 为光子角波数（$k = q = E_\gamma/\hbar c$），严格与电荷型算符 $\hat{\mathcal{M}}_S$ 配对：
\begin{equation}\label{eq:3-4-delta-def}
\delta(E2/M1) = \frac{\sqrt{3}}{10} \frac{E_\gamma}{\hbar c} \frac{\langle I_f \| \hat{\mathcal{M}}(E2) \| I_i \rangle}{\langle I_f \| \hat{\mathcal{M}}_S(M1) \| I_i \rangle} = \frac{\sqrt{3}}{10} \frac{E_\gamma}{\hbar} \frac{\langle I_f \| \hat{\mathcal{M}}(E2) \| I_i \rangle}{\langle I_f \| \hat{\mathcal{M}}_\mu(M1) \| I_i \rangle}.
\end{equation}
采用磁矩型算符 $\mathcal{M}_\mu$ 时，公式中等价提取为 $E_\gamma/\hbar$ 因子；以各自命名单位归一化的无量纲数值 $b_{M1}$ 完全相同，绝不应在量纲换算中再次额外施加 $c^2$ 因子；
\end{itemize}
\item \textbf{能级寿命测量与绝对跃迁几率（$B(E2)$ 和 $B(M1)$）}：利用多普勒位移衰减法（DSAM）或反冲距离法（RDM）测出激发态平均寿命 $\tau$，结合光子分枝比、内转换电子系数及严格反演的多极混合比 $\delta$，提取绝对跃迁几率 $B(E2)_{\text{inter}}$ 与 $B(M1)_{\text{inter}}$。在理想的纯转子与冻结顺排小振幅摆动模型的大自旋极限下（Bohr 与 Mottelson 1975，Frauendorf 与 D{\"o}nau 2014）~\cite{book-bohr-mottelson-nuclear-structure-vol2,frauendorf-2014-transverse-wobbling}，带间电四极跃迁表现出增强的集体强度，并在固定参数与领先阶渐近展开下呈现 $B(E2)_{\text{inter}} \sim 1/I$ 的标度趋势。然而必须清醒认识到，$1/I$ 标度仅是特定固定参数局部模型在大自旋渐近区的极限性质，并不普遍适用于所有自旋区间；同时，带间 $M1$ 跃迁强度的行为敏感依赖于磁偶极算符结构、单粒子与集体回转磁比差（$g_K - g_R$）以及微观核芯极化，Frauendorf 与 D{\"o}nau（2014）的 QTR 计算本身即表明 $M1$ 强度存在复杂的结构依赖，绝不能先验地将极弱 $M1$ 跃迁作为所有摆动模型的普适必然特征；
\item \textbf{微观与有效唯象理论模型的前向联合预言检验}：最终的物理裁决必须依靠自洽微观密度泛函、推转全三维准粒子转子模型、相互作用玻色子—费米子模型（IBFM）或微观投影展开等各类自洽唯象与微观有效理论模型，在统一且自洽的参数输入下，对能谱结构、衰变链分支、混合比以及电磁跃迁强度给出具备可证伪性的前向联合预言，并与竞争性物理机制（旋称分裂、组态交叉、软核涨落）进行严密的定量对账。不应将物理判决局限于某单一列举的微观计算框架。
\end{enumerate}

对非轴对称核转动态的物理确认，本质上是在不断提升的多维实验精密度与自洽理论约束下，对竞争性物理图像进行的严密证伪与综合鉴别过程。
